\chapter{原子的能级和辐射}
\problem{2.1}{试计算氢原子的第一玻尔轨道上电子绕核运动的频率、线速度和加速度。}
\answerarea
\problem{2.2}{试由氢原子的里德伯常数计算基态氢原子的电离电势和第一激发电势。}
\answerarea\clearpage
\problem{2.3}{用能量为 12.5 电子伏特的电子去激发基态氢原子。问受激发的氢原子向低能级跃迁时，会出现哪些波长的光谱线？}
\answerarea
\problem{2.4}{试估算一次电离的氦离子 $\mathrm{He}^+$、二次电离的锂离子 $\mathrm{Li}^{++}$ 的第一玻尔轨道半径、电离电势、第一激发电势和赖曼系第一条谱线波长分别与氢原子的上述物理量之比值。}
\answerarea\clearpage
\problem{2.5}{试问二次电离的锂离子 $\mathrm{Li}^{++}$ 从其第一激发态向基态跃迁时发出的光子，是否有可能使处于基态的一次电离的氦离子 $\mathrm{He}^+$ 中的电子电离掉？}
\answerarea
\problem{2.6}{氢与其同位素氘（质量数为 2）混在同一放电管中，摄下两种原子的光谱线。试问其巴尔末系的第一条（$\mathrm H_\alpha$）光谱线之间的波长差 $\Delta\lambda$ 有多大？已知氢的里德伯常数 $R_{\mathrm H}=1.0967758\times10^7\,\mathrm{m^{-1}}$，氘的里德伯常数 $R_{\mathrm D}=1.0970742\times10^7\,\mathrm{m^{-1}}$。}
\answerarea\clearpage
\problem{2.7}{已知一对正负电子绕其共同的质心转动会暂时形成类似于氢原子结构的“正电子素”。试计算“正电子素”由第一激发态向基态跃迁发射光谱的波长 $\lambda$ 为多少 $\angstrom$？}
\answerarea
\problem{2.8}{试证明氢原子中的电子从 $n+1$ 轨道跃迁到 $n$ 轨道，发射光子的频率 $\nu_n$，当 $n\gg1$ 时光子频率即为电子绕第 $n$ 玻尔轨道转动的频率。}
\answerarea\clearpage
\problem{2.9}{$\mathrm{Li}$ 原子序数 $Z=3$，其光谱的主线系可用下式表示：
\[
\widetilde\nu=\frac{R}{(1+0.5951)^2}-\frac{R}{(n-0.0401)^2}.
\]
已知 $\mathrm{Li}$ 原子电离成 $\mathrm{Li}^{++}$ 离子需要 203.44 电子伏特的功。问如要把 $\mathrm{Li}^{+}$ 离子电离为 $\mathrm{Li}^{++}$ 离子，需要多少电子伏特的功？}
\answerarea
\problem{2.10}{具有磁矩的原子，在横向均匀磁场和横向非均匀磁场中运动时有什么不同？}
\answerarea\clearpage
\problem{2.11}{在史特恩--革拉赫实验中，处于基态的窄银原子束通过不均匀横向磁场，磁场梯度为 $\partial B/\partial z=10^3\,\mathrm{T/m}$，磁极纵向范围 $L_1=0.04\,\mathrm m$（习题图 2.1），从磁极到屏距离 $L_2=0.10\,\mathrm m$，原子的速度 $v=5\times10^2\,\mathrm{m/s}$。在屏上两束分开的距离 $d=0.002\,\mathrm m$，试确定原子磁矩在磁场方向上投影 $\mu$ 的大小（设磁场边缘的影响可忽略不计）。
\fig{2_1}{2.1}}
\answerarea
\problem{2.12}{观察真空玻璃管中由激发原子束所发光谱线的强度沿原子射线束的减弱情况，可以测定各激发态的平均寿命。若已知原子束中原子速度 $v=10^3\,\mathrm{m/s}$，在沿粒子束方向上相距 1.5 毫米共振光谱线强度减少到 $1/3.32$，试计算这种原子在共振激发态的平均寿命。}
\answerarea\clearpage
